% !PRINTED 384
\weberpage{384}
\providecommand{\weberfourup}{\mathord{\tikz[baseline=-0.2ex,x=0.23em,y=0.23em,line width=0.45pt]{\draw (0,0)--(-1.05,1) (0,0)--(-0.35,1.08) (0,0)--(0.35,1.08) (0,0)--(1.05,1);}}}
\begin{center}
\S.\,102.

\gesperrt{Darstellung der Gruppe.}
\end{center}

Der Grad der Gruppe $P$ ist sofort anzugeben. Da es 288 vollständige Siebener-Systeme giebt, und da die Elemente eines solchen Systemes auf $\Pi(7)$ Arten permutirt werden können, so ist der Grad der Gruppe:

\[
288\,\Pi(7) = 36\,\Pi(8) = 1\,451\,520.
\]

Bei der Bildung der Permutationen der Wurzeln benutzen wir zweierlei Bezeichnung. Zunächst

\weq{1}{\xi_1,\ \xi_2,\ \dots,\ \xi_7,\ \xi_{ik},}

worin $i,k$ von 1 bis 7 geht, und die einheitliche Bezeichnung $[i\,k]$, wobei $i,k$ von 1 bis 8 geht, und wobei $\xi_1,\xi_2\dots$ durch $[1\,8],[2\,8],\dots$ zu bezeichnen sind.

Wenn wir nun zwei Ziffern aus der Reihe $1,2,\dots,7$ permutiren, z.\,B. 1 mit 2, so geht $\xi_1$ in $\xi_2$ über, und nach der Definition \S.\,99 bleibt $\xi_{12}$ ungeändert, $\xi_{13}$ geht in $\xi_{23}$ über u.\,s.\,f.

Wenn wir also in der Reihe der Wurzeln $[i\,k]$ die Ziffern 1 bis 7 beliebig permutiren, so erhalten wir lauter Permutationen der Gruppe $P$.

Nun haben wir im \S.\,100 gesehen, dass bei der Anordnung in syzygetische und azygetische Tripel und in Folge dessen auch in Complexe und vollständige Systeme die Ziffer 8 mit den übrigen Ziffern 1 bis 7 vollständig gleichberechtigt auftritt, und da die Zuordnung der Wurzel $\xi_{ik}$ zu dem Paare $\xi_i\xi_k$ [durch die Formel \S.\,101, (5)] nach \S.\,99 nur von dieser Anordnung abhängt, so können wir auch die Ziffern $1,2,\dots,7,8$ permutiren, ohne dass wir aus der Gruppe $P$ herauskommen. Es ist also in $P$ ein Theiler enthalten, der mit der symmetrischen Gruppe der Permutationen von acht Elementen isomorph ist, der also den Index 36 hat und den wir mit $S$ bezeichnen wollen.

Um die noch fehlenden Permutationen von $P$ zu bestimmen, lassen wir an Stelle des Siebener-Systemes $\xi_1,\dots,\xi_7$ ein neues vom Typus $\Delta\,\weberfourup$ treten, etwa so:

% [sic] The source prints a full stop after [6 8] in the first row.
\weq{2}{\begin{pmatrix} [1\,8], & [2\,8], & [3\,8], & [4\,8], & [5\,8], & [6\,8]. & [7\,8] \\ {}[2\,3], & [3\,1], & [1\,2], & [4\,8], & [5\,8], & [6\,8], & [7\,8] \end{pmatrix},}
%CONT>>
